% Provenance (never printed): sources/book2/BRIEFS.md, g11-08. Hidden quality (the lemons model), hidden action (moral
% hazard, principal and agent), answers on each side, compulsory and voluntary cover. Real numbers:
% figdata/grade-11/08-asymmetric-information/ (RAND experiment, cited from Aron-Dine, Einav and Finkelstein 2012, Table 2;
% a seeded lemons simulation, invented). Ledger: HIE design (Wikipedia, RAND R-3476 abstract, AEF 2012), Bond 1982. No
% certification or warranty study was read, so none is cited; the signalling diagram is replaced by a worked example.
% All Brenland numbers are invented.
\chapter{Asymmetric Information}\label{ch:g11:asymmetric-information}

From 1974 to 1982, several thousand Americans took part in an experiment in which their health insurance was chosen by lot. Some paid nothing for their care; others
paid 25, 50 or 95 per cent of each bill, up to a yearly limit. Those whose care was free spent 2170 dollars a year each (in 2011 dollars); those who paid 95 per cent
spent 1325, 39 per cent less. Because the plans were drawn by lot, the groups were alike in health on average; what differed was what they
paid at the counter.
% ledger: g11:hie-design, g11:hie-free, g11:hie-95; figdata: hie (free 2170, c95 1325)

Insurance changes what people do, and an insurer cannot watch them. A seller knows more about the car he sells than the buyer does. When one side of a deal knows
something the other does not, markets can work badly, or not at all. This chapter shows how, and the answers that markets and governments have found.

\begin{openingchart}{The RAND Health Insurance Experiment: mean annual medical spending per person (2011 dollars) and share of people with any medical spending in the year,
by the share of each bill the family paid. Each plan's figure is the free plan's mean plus the estimated difference from it; the gap between the 25\,\% and 50\,\% plans is
within the estimates' margin of error.}{Aron-Dine, Einav and Finkelstein (2012), Table 2, from the RAND Health Insurance Experiment}
\begin{tikzpicture}
\begin{groupplot}[group style={group size=2 by 1, horizontal sep=1.6cm}, width=0.47\linewidth, height=5.0cm, ybar, /pgf/bar width=14pt,
  xtick={0,1,2,3}, xticklabels={free, 25\,\%, 50\,\%, 95\,\%}, enlarge x limits=0.2, ymin=0, xlabel={share of the bill paid},
  nodes near coords, every node near coord/.append style={font=\scriptsize}, x tick label style={/tikz/font=\scriptsize}]
\nextgroupplot[width=0.47\linewidth, ymax=2600, ylabel={dollars a year}, title={\small spending per person},
  nodes near coords={\pgfmathprintnumber[fixed, precision=0, 1000 sep={}]{\pgfplotspointmeta}}, scaled y ticks=false,
  y tick label style={/pgf/number format/1000 sep={}}]
\addplot[fill=omDef!60, draw=omDef] table[x=x, y=spend] {figdata/grade-11/08-asymmetric-information/out/hie.csv};
\nextgroupplot[width=0.47\linewidth, ymax=110, ylabel={per cent}, title={\small with any spending}]
\addplot[fill=omThm!60, draw=omThm] table[x=x, y=anyuse_pct] {figdata/grade-11/08-asymmetric-information/out/hie.csv};
\end{groupplot}
\end{tikzpicture}
\end{openingchart}

\begin{recall}
From the grades 6--9 volume: a market brings together buyers and sellers who agree on a price. From \cref{ch:g10:surplus-and-efficiency}: \omterm{def:g10:surplus-and-efficiency:wtp}{willingness to pay}, and
the surplus lost when trades worth making are not made.
\end{recall}

\begin{definition}[Asymmetric information]\label{def:g11:asymmetric-information:asymmetric-information}
There is \emph{asymmetric information}\index{asymmetric information} in a deal when one party knows something relevant to it that the other does not.
\end{definition}

\section{Hidden quality}\label{sec:g11:asymmetric-information:lemons}

\begin{definition}[Adverse selection]\label{def:g11:asymmetric-information:adverse-selection}
\emph{Adverse selection}\index{adverse selection} occurs when the hidden characteristics of those who choose to trade make the deal worse for the uninformed side:
the cars offered for sale are the worse ones, the people who buy insurance are the riskier ones.
\end{definition}

\begin{example}[Used cars in Brenland]\label{ex:g11:asymmetric-information:cars}
The qualities of a town's used cars are spread evenly between 2000 and 10\,000 dalens; a seller values his car at its quality $q$, a buyer at $1.5q$. Every car is worth
more to a buyer than to its seller, so all would be sold if buyers could see quality. They cannot. At a price $p$, only the owners of cars with $q \le p$ sell; the
average quality offered is $(2000 + p)/2$, which a buyer values at $1.5 \times (2000 + p)/2 = 1500 + 0.75p$. Buyers pay $p$ only if that is at least $p$, that is
$p \le 6000$. At 6000, the cars between 6000 and 10\,000 stay with their owners: half the cars, and the better half, are never sold.
\end{example}

\begin{omfigure}
\begin{tikzpicture}[econ, x=0.0009cm, y=0.0005cm]
\draw[econ axes] (2000,10500) node[left, font=\small] {dalens} |- (11000,2000) node[below left=12pt and 0pt, font=\small] {price $p$};
\draw[supply] (2000,2000) -- (10000,10000) node[right, font=\scriptsize] {price $p$};
\draw[demand] (2000,3000) -- (10000,9000) node[right, font=\scriptsize] {buyers' value of the average car offered};
\node[eqpoint] at (6000,6000) {}; \draw[guide] (6000,2000) -- (6000,6000) -- (2000,6000);
\node[below, font=\footnotesize] at (2000,2000) {2000}; \node[below, font=\footnotesize] at (6000,2000) {6000}; \node[below, font=\footnotesize] at (10000,2000) {10\,000};
\foreach \y in {3000,6000,9000} \draw (2000,\y) -- (1900,\y) node[left, font=\footnotesize] {\y};
\end{tikzpicture}

{\small The Brenland market. Above a price of 6000, buyers value the average car offered at less than its price, and they do not buy; the price settles at 6000.}
\end{omfigure}

\begin{model}{The lemons model}
\assumes Sellers know their good's quality; buyers know only the average quality of what is offered; no one can check quality before buying.
\predicts The best goods are withdrawn, the price falls with the average quality, and the market may shrink to the worst goods or vanish.
\tested Little support in the one test cited: in a 1977 survey of American pickup trucks, trucks bought used needed no more maintenance than trucks of
  the same age and mileage that had never changed hands (Bond 1982). The existence of inspections, warranties and dealers' guarantees is consistent with the
  problem biting elsewhere, since sellers and buyers pay to get round it.
\breaks When quality can be checked (inspection, a known history), when sellers have a reputation to lose, or when they offer a warranty.
\end{model}
% ledger: g11:bond1982

\begin{omchart}{An invented used-car market, simulated: 1000 cars with qualities drawn at random (fixed seed), sellers valuing a car at its quality, buyers at 1.5 times
it. Each round, buyers offer 1.5 times the average quality of the cars offered at the last round's price. With qualities between 2000 and 10\,000 dalens the price
settles near 6000; with qualities between 0 and 10\,000 the market unravels.}{invented data, seeded simulation}
\begin{tikzpicture}
\begin{axis}[omtime, width=0.84\linewidth, height=5.2cm, xmin=0, xmax=15, ymin=0, ymax=10000, xlabel={round}, ylabel={price (dalens)}, scaled y ticks=false,
  y tick label style={/pgf/number format/1000 sep={\,}}]
\addplot[omDef, thick, mark=*, mark size=1pt] table[x=round, y=floor] {figdata/grade-11/08-asymmetric-information/out/lemons.csv}
  node[pos=0.8, above, font=\scriptsize] {qualities 2000 to 10\,000};
\addplot[omThm, thick, mark=*, mark size=1pt] table[x=round, y=nofloor] {figdata/grade-11/08-asymmetric-information/out/lemons.csv}
  node[pos=0.4, anchor=south west, font=\scriptsize] {qualities 0 to 10\,000};
\end{axis}
\end{tikzpicture}
\end{omchart}
% figdata: lemons, lemons_end (5927, 0.514; 221, 0.02)

\section{Hidden action}\label{sec:g11:asymmetric-information:moral-hazard}

\begin{definition}[Moral hazard]\label{def:g11:asymmetric-information:moral-hazard}
\emph{Moral hazard}\index{moral hazard} occurs when a party whose actions cannot be observed takes less care, or more risk, because someone else bears part of the cost.
\end{definition}

\begin{example}[Care and insurance]\label{ex:g11:asymmetric-information:care}
A cyclist's bicycle, worth 10\,000 dalens, is stolen with probability 10\,\% a year if he leaves it unlocked and 4\,\% if he always locks it, which costs him trouble worth
300 dalens a year. Uninsured, locking cuts his expected loss by $0.06 \times 10\,000 = 600 > 300$: he locks. Fully insured, locking saves him nothing: he does not lock,
and the insurer's expected payout rises from 400 to 1000 dalens a year. If he must pay the first $D$ dalens of any loss, locking saves him $0.06D$, and he locks when
$0.06D \ge 300$, that is when $D \ge 5000$.
\end{example}

\begin{definition}[Deductible]\label{def:g11:asymmetric-information:deductible}
A \emph{deductible}\index{deductible} is the part of each loss that the insured person pays before the insurance pays the rest.
\end{definition}

\begin{omfigure}
\begin{tikzpicture}[econ, x=0.0008cm, y=0.0055cm]
\draw[econ axes] (0,700) node[left, font=\small] {dalens a year} |- (11000,0) node[below left=12pt and 0pt, font=\small] {deductible $D$};
\draw[demand] (0,0) -- (10000,600) node[right, font=\scriptsize] {saving from locking, $0.06D$};
\draw[supply] (0,300) -- (10000,300) node[right, font=\scriptsize] {cost of locking};
\node[eqpoint] at (5000,300) {}; \draw[guide] (5000,0) -- (5000,300);
\node[below, font=\footnotesize] at (5000,0) {5000}; \node[below, font=\footnotesize] at (10000,0) {10\,000};
\foreach \y in {300,600} \draw (0,\y) -- (-100,\y) node[left, font=\footnotesize] {\y};
\node[font=\scriptsize, align=center] at (2600,235) {no care}; \node[font=\scriptsize, align=center] at (7700,385) {care};
\end{tikzpicture}

{\small The cyclist's choice. Below a \omterm{def:g11:asymmetric-information:deductible}{deductible} of 5000 dalens the saving from locking is less than its cost, and he does not lock.}
\end{omfigure}

\begin{model}{Moral hazard in insurance}
\assumes The insurer cannot observe the insured person's care or use; the insured person weighs only what he himself pays.
\predicts The more of the cost the insurance bears, the less care and the more use.
\tested In the RAND experiment, families assigned at random to plans in which they paid a larger share of each bill spent less: 2170 dollars a year with free care,
  1325 when they paid 95\,\% (up to the yearly limit), and fewer of them used any care at all.
\breaks When the insurer can observe care (a lock it inspects, a fitted alarm), or when the use is not chosen (an emergency operation).
\end{model}
% ledger: g11:hie-free, g11:hie-95; figdata: hie

\begin{record}[The RAND Health Insurance Experiment, 1974--82]
% ledger: g11:hie-design, g11:hie-free, g11:hie-25, g11:hie-50, g11:hie-95, b2:randhie2006; figdata: hie
Families were assigned by lot to plans in which they paid 0, 25, 50 or 95\,\% of each bill, with their out-of-pocket spending capped at a share of family income, never more than 1000 dollars of the time, and were
covered for three to five years; 5809 people were in these plans. Spending was highest with free care and lowest with the 95\,\% plan, and the largest drop came
between free care and the first plan with any charge. Because the plans were drawn by lot, the difference measures the effect of the price paid at the counter, not
the different health of those who choose generous insurance.
\end{record}

\begin{definition}[Principal--agent problem]\label{def:g11:asymmetric-information:principal-agent}
A \emph{principal--agent problem}\index{principal--agent problem} arises when one party, the agent, acts on behalf of another, the principal, who cannot fully observe
the agent's effort, and their interests differ: a manager and the owners of a firm, a builder and his client.
\end{definition}

Owners can tie a manager's pay to the firm's results, so that the manager gains when they do. The fix has a price: results depend on luck as well as effort, so the
manager bears risk and asks to be paid for it, and pay tied to a measure invites effort on the measure rather than on the firm.

\section{Answers}\label{sec:g11:asymmetric-information:answers}

\begin{definition}[Signalling and screening]\label{def:g11:asymmetric-information:signalling}
\emph{Signalling}\index{signalling} is an action by the informed party, costly enough that only those of the good type find it worth taking, which reveals their
type. \emph{Screening}\index{screening}\phantomsection\label{def:g11:asymmetric-information:screening} is a choice offered by the uninformed party, designed so
that different types choose differently and so reveal themselves.
\end{definition}

\begin{example}[A signal that separates]\label{ex:g11:asymmetric-information:signal}
Workers are of two kinds, producing 30 or 20 thousand dalens a year; employers cannot tell them apart. A certificate costs the able 8 thousand in effort and the others 15.
If employers pay 30 to holders and 20 to the rest, an able worker gains $30 - 8 = 22 > 20$ and takes it; another would get $30 - 15 = 15 < 20$ and does not. The
certificate need not raise anyone's output to separate them: what matters is that it costs the less able more.
\end{example}

The informed side can also offer a warranty (only a seller who knows his car is good can afford to promise repairs), build a reputation, or pay for an inspection that
certifies quality. The uninformed side screens: an insurer offers a cheap policy with a large \omterm{def:g11:asymmetric-information:deductible}{deductible} and a dear one with a small \omterm{def:g11:asymmetric-information:deductible}{deductible}, and careful drivers
choose the first; it also prices by experience, raising the premiums of those who have claimed.

\section{Compulsory and voluntary cover}\label{sec:g11:asymmetric-information:compulsory}

\omterm{def:g11:asymmetric-information:adverse-selection}{Adverse selection} hits insurance hard. If one premium is charged to everyone and buying is voluntary, the low risks find it too dear and leave; the premium must rise to
cover those who stay, and more leave. A pool can shrink to the highest risks. Compulsory cover at one premium, such as a health scheme that every resident must join and that
charges everyone the same, keeps the low risks in the pool and the premium near the average. Its cost falls on the low risks, who pay more than
their expected loss. (Compulsory motor insurance is another matter: it is priced by risk, and it is compelled mainly so that the people a driver
harms are paid.) Voluntary insurance survives
where insurers can measure risk and price it, by age, by record, by the car, and where people's private knowledge of their own risk is small.

\section{Exercises}

\begin{exercise}[$\star$]\label{exo:g11:asymmetric-information:1}
\omterm{def:g11:asymmetric-information:adverse-selection}{Adverse selection} or \omterm{def:g11:asymmetric-information:moral-hazard}{moral hazard}? A buyer of life insurance hides a serious illness; a cyclist with theft insurance leaves his bicycle unlocked; a seller of a used phone
knows its battery is failing.
\end{exercise}

\begin{exercise}[$\star$]\label{exo:g11:asymmetric-information:2}
\omterm{def:g11:asymmetric-information:adverse-selection}{Adverse selection} or \omterm{def:g11:asymmetric-information:moral-hazard}{moral hazard}? A worker slacks when his manager is away; the people who buy dental insurance are mostly those who expect to need a dentist; a bank whose
deposits are insured by the state makes riskier loans.
\end{exercise}

\begin{exercise}[$\star$]\label{exo:g11:asymmetric-information:3}
\omterm{def:g11:asymmetric-information:adverse-selection}{Adverse selection} or \omterm{def:g11:asymmetric-information:moral-hazard}{moral hazard}? A driver of a hired car drives it harder than his own; a borrower whose plan is risky is keener to borrow at a high interest rate than one
whose plan is safe; a homeowner with fire insurance stops checking his wiring.
\end{exercise}

\begin{exercise}[$\star$]\label{exo:g11:asymmetric-information:4}
From the opening chart, compute in per cent how much less a person spent on the 95\,\% plan than on the free plan.
\end{exercise}

\begin{exercise}[$\star$]\label{exo:g11:asymmetric-information:5}
From the opening chart, by how many points does the share of people with any medical spending fall between the free and the 95\,\% plan? Why should one not conclude that
the 50\,\% plan made people spend more than the 25\,\% plan?
\end{exercise}

\begin{exercise}[$\star\star$]\label{exo:g11:asymmetric-information:6}
Half the used cars of a town are good, worth 8000 dalens to their sellers and 10\,000 to buyers; half are bad, worth 2000 to sellers and 3000 to buyers. Buyers cannot tell
them apart. Which cars are sold, and at what price?
\end{exercise}

\begin{exercise}[$\star\star$]\label{exo:g11:asymmetric-information:7}
Same question if 80\,\% of the cars are good.
\end{exercise}

\begin{exercise}[$\star\star$]\label{exo:g11:asymmetric-information:8}
In \cref{exo:g11:asymmetric-information:6}, a seller may offer a one-year warranty. It would cost the seller of a good car 200 dalens in expected repairs and the seller of a
bad car 8000. Show that if buyers pay 10\,000 for a car with a warranty, only good cars carry one.
\end{exercise}

\begin{exercise}[$\star\star$]\label{exo:g11:asymmetric-information:9}
A loss of 2000 dalens strikes with probability 10\,\% a year. A policy has a \omterm{def:g11:asymmetric-information:deductible}{deductible} of 500. Compute the insurer's expected payout and the insured person's expected
out-of-pocket cost.
\end{exercise}

\begin{exercise}[$\star\star$]\label{exo:g11:asymmetric-information:10}
In the certificate example, let the certificate cost the able worker $c_H$ and the other $c_L$ (thousand dalens). For which values does the certificate separate the two
kinds at wages of 30 and 20?
\end{exercise}

\begin{exercise}[$\star\star$]\label{exo:g11:asymmetric-information:11}
From the simulated chart, read the price in round 0 and round 15 for each market. In the market whose qualities start at 2000, what share of the cars is traded at the
end? Compare with the example of the chapter.
\end{exercise}

\begin{exercise}[$\star\star\star$]\label{exo:g11:asymmetric-information:12}
What does the lemons model predict about the maintenance of trucks bought used compared with trucks that never changed hands? What did Bond (1982) find? Give two
reasons the used-truck market might escape the lemons problem.
\end{exercise}

\begin{exercise}[$\star\star\star$]\label{exo:g11:asymmetric-information:13}
Three equal groups of people have expected losses of 100, 300 and 800 dalens a year; each will pay up to 50 dalens more than its expected loss for full insurance. An
insurer charges everyone the average expected loss of those it insures. Show that the pool shrinks to the highest risks.
\end{exercise}

\begin{exercise}[$\star\star\star$]\label{exo:g11:asymmetric-information:14}
A manager's effort raises a firm's expected profit from 1000 to 1600 thousand dalens a year and costs him the equivalent of 200 thousand. On a fixed salary, does he make
the effort? What share $s$ of profit must his pay include for him to make it? What does this cost the owners, and what new problem does it create?
\end{exercise}

\begin{exercise}[$\star\star\star$]\label{exo:g11:asymmetric-information:15}
In \cref{exo:g11:asymmetric-information:13}, insurance is made compulsory at a premium equal to the average expected loss of all three groups. Find the premium, and say
who gains and who loses compared with the voluntary market.
\end{exercise}

\section{Problem: Used Cars in a Town in Brenland}

\begin{problem}[{Weekend problem --- hidden quality and an inspection service}]
\label{pb:g11:asymmetric-information:1}
A town in Brenland has many used cars, with qualities $q$ spread evenly between 2000 and 10\,000 dalens. A seller values a car at $q$, a buyer at $1.5q$, and there are more
buyers than cars.

\medskip
\noindent\textbf{Part I --- Hidden quality.}
\begin{enumerate}
\item If buyers could see quality, which cars would be sold?
\item What would be the average gain from trade per car?
\item Buyers cannot see quality. At a price $p$ between 2000 and 10\,000, which cars are offered?
\item What is the average quality of the cars offered at price $p$?
\item How much is the average car offered worth to a buyer?
\end{enumerate}

\medskip
\noindent\textbf{Part II --- Where the market settles.}
\begin{enumerate}[resume]
\item Find the highest price buyers are willing to pay.
\item What share of the cars is traded?
\item What is the average quality of the cars traded, and the buyers' average gain?
\item Compute the gain from trade realised, per car in the town, and compare it with question 2.
\item If qualities were spread between 0 and 10\,000, where would the price settle?
\end{enumerate}

\medskip
\noindent\textbf{Part III --- An inspection service.}
\begin{enumerate}[resume]
\item An inspection, paid by the seller, costs 400 dalens and shows the car's quality to all; buyers then pay $1.5q$. Which sellers inspect?
\item What share of the cars is traded, and what is the average price?
\item What is the gain from trade realised per car, net of the inspection?
\item By how much does the service raise the gain per car?
\item Why might a warranty do the same work as the inspection?
\end{enumerate}

\medskip
\noindent\textbf{Part IV --- Variations.}
\begin{enumerate}[resume]
\item Without inspection, buyers value a car at $1.2q$. Find the price and the share traded.
\item Without inspection, buyers value a car at $2q$. Show that all cars are traded.
\item Explain the difference between the two answers.
\item A dealer sells hundreds of cars a year in the town. Why may he not sell bad cars as good ones, even without inspection?
\item State the result of the problem: the price at which the market settles, and the share of cars traded, with and without inspection.
\end{enumerate}
\end{problem}
